% Currículo de Cheila Betina Schilling dos Santos
% Primeira versão editorial — conteúdo sujeito a revisão antes da publicação.
\documentclass[10pt,a4paper]{article}

\usepackage[T1]{fontenc}
\usepackage[utf8]{inputenc}
\usepackage[brazilian]{babel}
\usepackage[a4paper,top=1.45cm,bottom=1.35cm,left=1.55cm,right=1.55cm]{geometry}
\usepackage{lmodern}
\usepackage{xcolor}
\usepackage{hyperref}
\usepackage{enumitem}
\usepackage{tabularx}
\usepackage{array}
\usepackage{microtype}
\usepackage{titlesec}
\usepackage{parskip}
\usepackage{needspace}
\usepackage{multicol}

\definecolor{ink}{HTML}{171612}
\definecolor{red}{HTML}{A43A2B}
\definecolor{grey}{HTML}{656057}
\definecolor{paper}{HTML}{F8F3E8}

\hypersetup{
  colorlinks=true,
  urlcolor=red,
  linkcolor=red,
  pdftitle={Currículo — Cheila Betina Schilling dos Santos},
  pdfauthor={Cheila Betina Schilling dos Santos}
}

\pagestyle{empty}
\setlength{\parindent}{0pt}
\setlist[itemize]{leftmargin=1.15em,itemsep=1pt,topsep=3pt,parsep=0pt}
\setlength{\columnsep}{18pt}
\setlength{\emergencystretch}{2em}
\raggedbottom
\titleformat{\section}{\large\bfseries\color{red}\sffamily}{}{0pt}{}[\vspace{-3pt}\color{red}\rule{\textwidth}{0.5pt}]
\titlespacing*{\section}{0pt}{10pt}{5pt}
\titleformat{\subsection}{\normalsize\bfseries\color{ink}\sffamily}{}{0pt}{}
\titlespacing*{\subsection}{0pt}{7pt}{1pt}

\newcommand{\role}[3]{%
  \Needspace{13\baselineskip}%
  \textbf{#1}\\[-1pt]
  {\small\textcolor{grey}{#2}}\\[-1pt]
  #3\par}
\newcommand{\skill}[2]{\textbf{#1:} #2\par}

\begin{document}

{\sffamily
{\fontsize{24}{27}\selectfont\bfseries Cheila Betina Schilling dos Santos}\par
\vspace{3pt}
{\large\color{red} Dados · Analytics · Advanced Analytics}\par
\vspace{5pt}
{\small
Araricá, RS \quad | \quad (+55) 51 99949-2331 \quad | \quad
\href{mailto:betina.ssc@gmail.com}{betina.ssc@gmail.com}\\
\href{https://github.com/betinaschilling}{github.com/betinaschilling} \quad | \quad
\href{https://betinaschilling.github.io/}{betinaschilling.github.io} \quad | \quad
\href{https://www.linkedin.com/in/cheila-betina-santos/}{linkedin.com/in/cheila-betina-santos} \quad | \quad
\href{https://orcid.org/0009-0008-6145-3587}{ORCID}
}\par
}

\section{Resumo profissional}

Profissional de dados com experiência na interseção entre engenharia, ciência de dados e analytics. Atua na definição de problemas, desenho de soluções, escolha metodológica, construção de pipelines e modelos, comunicação de resultados e evolução de produtos analíticos. Combina conhecimento técnico em dados distribuídos, cloud, machine learning, forecasting e advanced analytics com experiência em decisões de negócio, governança, qualidade, implantação e colaboração multidisciplinar.

\section{Experiência profissional}

\role{Lojas Renner S.A. — Cientista de Dados Pleno}{Porto Alegre, RS · ago. 2025 – atual}{
\begin{itemize}
  \item Desenvolvimento de modelos preditivos para forecasting, elasticidade, pricing e alocação, além de modelos de classificação aplicados a churn.
  \item Aplicação de métodos estatísticos e de séries temporais, incluindo regressão, OLS, ARIMA, SARIMA, Prophet e VAR/VECM.
  \item Utilização de CatBoost, XGBoost e LightGBM, além de Chronos, TimesFM e DeepAR, tanto em experimentos quanto em soluções utilizadas em produção.
  \item Construção de códigos modularizados em Python e Spark para ingestão, processamento, treinamento e avaliação de modelos.
  \item Integração de dados provenientes de Databricks, APIs públicas e fontes externas, incluindo variáveis climáticas e indicadores econômicos.
  \item Execução de análises exploratórias e estatísticas, engenharia de atributos, seleção de variáveis, ajuste de hiperparâmetros e avaliação de modelos.
  \item Participação na produtização, implantação e acompanhamento de modelos em produção.
  \item Acompanhamento e sustentação de produtos de dados novos e legados, incluindo monitoramento de pipelines, modelos e rotinas operacionais.
  \item Integração de pipelines com mecanismos de monitoramento e alertas, incluindo Zabbix, para identificação de falhas de execução.
\end{itemize}}

\role{Grendene S.A. — Analista de Dados Sênior (Produto) / Engenheira de Dados (COE)}{Farroupilha, RS · abr. 2024 – ago. 2025}{
\begin{itemize}
  \item Desenvolvimento e otimização de pipelines de dados utilizando Informatica Data Integration, Apache Spark, AlloyDB e Google Cloud Platform.
  \item Implementação de soluções de processamento em streaming com Pub/Sub e arquiteturas de self-service SQL para dados batch e tempo real via Analytics Hub.
  \item Estruturação e gestão de camadas de dados — stage, raw, silver, gold, self-service e agregadas — utilizando Cloud Storage.
  \item Desenvolvimento e manutenção de fluxos de dados para integração entre fontes operacionais, camadas analíticas e produtos de dados.
  \item Implantação e manutenção de painéis Power BI em ambientes de homologação e produção.
  \item Participação em comitês de aprovação de mudanças de código e deploys em produção, incluindo revisão de merge requests no GitLab.
  \item Treinamento de equipes analíticas em boas práticas de desenvolvimento, qualidade de código e versionamento.
  \item Aplicação de conceitos de IAM para governança, segurança e controle de acessos aos ambientes e dados.
\end{itemize}}

\role{Portobello Shop — Analista de Dados Pleno (COE) / Analytics Engineer (Produto)}{Florianópolis, SC · ago. 2022 – abr. 2024}{
\begin{itemize}
  \item Desenvolvimento de dashboards interativos em Tableau para análise de vendas e acompanhamento de desempenho operacional.
  \item Execução de análises exploratórias sobre grandes volumes de dados utilizando Databricks.
  \item Desenvolvimento e manutenção de consultas SQL nas camadas gold e agregadas, com implementação de regras de negócio.
  \item Administração de usuários, grupos e permissões no Tableau Server e implantação de painéis em produção.
  \item Condução de treinamentos para equipes de negócio sobre dashboards autossustentáveis e fontes de dados integradas.
  \item Implementação de soluções de integração e automação de dados utilizando Apache Airflow.
\end{itemize}}

\role{Azzas S.A. — Analista de Dados / Estatística (Produto)}{Campo Bom, RS · dez. 2017 – ago. 2022}{
\begin{itemize}
  \item Desenvolvimento de modelos estatísticos para orçamento de receita, análise de variações de desempenho e revisão de projeções financeiras.
  \item Execução de análises de previsto versus realizado para apoiar ajustes nas estratégias comerciais e de vendas.
  \item Desenvolvimento de modelos de séries temporais com Prophet para previsão de volumetria mensal de pedidos.
  \item Criação de dashboards Tableau para planejamento, comercial, e-commerce e gestão.
  \item Validação de dados e atuação como data owner dos dados de omnichannel, apoiando governança e consistência entre canais físicos e digitais.
  \item Desenvolvimento e evolução de soluções analíticas para produtos de omnichannel e CRM, incluindo indicadores de NPS.
  \item Aplicação de regressão linear, classificação e análise de clusters para identificar padrões de comportamento e apoiar estratégias comerciais.
\end{itemize}}

\role{GetNet Tecnologia — Analista de Dados / Planejamento Operacional}{Porto Alegre, RS · fev. 2010 – ago. 2015}{
\begin{itemize}
  \item Desenvolvimento de modelos de previsão financeira e operacional para análise de fluxo de caixa, avaliação de riscos e planejamento de capacidade em processos de pagamentos.
  \item Aplicação de conceitos de pesquisa operacional e teoria das filas, incluindo Erlang-B e Erlang-C, para estimar demanda, dimensionar recursos e otimizar a capacidade de atendimento em centrais de chamadas.
  \item Criação de relatórios executivos e operacionais utilizando Power BI, SQL Server, VBA e automações com VB e RPA.
  \item Implantação de soluções de BI em SAP para análise financeira em D-1 e desenvolvimento de consultas SQL em Oracle.
  \item Utilização de técnicas de regressão para análise de tendências financeiras e de mercado.
\end{itemize}}

\section{Educação}

\begin{tabularx}{\textwidth}{@{}Xr@{}}
\textbf{Mestrado em Administração — foco em Ciência de Dados} & mar. 2026 – mar. 2029\\
FEEVALE · Novo Hamburgo, RS · em andamento &\\[3pt]
\textbf{Bacharelado em Estatística — segunda graduação} & mar. 2026 – mar. 2029\\
UNICV · Curitiba, PR · em andamento &\\[3pt]
\textbf{Bacharelado em Engenharia de Produção — segunda graduação} & mar. 2026 – mar. 2029\\
UNICV · Curitiba, PR · em andamento &\\[3pt]
\textbf{Pós-graduação em Matemática Aplicada} & jan. 2025 – maio 2025\\
Focus · São Paulo, SP · concluída &\\[3pt]
\textbf{Pós-graduação em Data Analytics} & abr. 2024 – abr. 2025\\
FIAP · São Paulo, SP · concluída &\\[3pt]
\textbf{Graduação em Processos Gerenciais} & jan. 2016 – jan. 2019\\
UNINTER · Curitiba, PR · concluída &
\end{tabularx}

\Needspace{18\baselineskip}
\section{Competências técnicas}

\begin{multicols}{2}
\raggedcolumns
\skill{Engenharia e plataformas de dados}{Python, SQL, Apache Spark, Databricks, BigQuery, Cloud Storage, AlloyDB, Pub/Sub, Analytics Hub, Oracle, PostgreSQL, MySQL e MongoDB.}
\skill{Cloud e operação}{Google Cloud Platform, AWS, Azure, Airflow, Cloud Run, Informatica Data Integration, pipelines batch e streaming, IAM, governança e controle de acessos.}
\skill{Ciência de dados e advanced analytics}{Regressão, classificação, clustering, forecasting, séries temporais, elasticidade, pricing, alocação, churn, otimização e análise de filas.}
\skill{Modelos estatísticos e preditivos}{OLS, ARIMA, SARIMA, Prophet, VAR/VECM, CatBoost, XGBoost, LightGBM, Chronos, TimesFM e DeepAR.}
\skill{Seleção e otimização}{Boruta, Optuna, engenharia de atributos, validação temporal, ajuste de hiperparâmetros e avaliação comparativa de modelos.}
\skill{Análise e visualização}{EDA, ADF, KPSS, CCF, Granger, cointegração, Tableau, Power BI, Matplotlib, Plotly e Streamlit.}
\skill{Métodos quantitativos}{Erlang-B, Erlang-C, teoria das filas, dimensionamento de capacidade e otimização de recursos.}
\skill{Engenharia de software e colaboração}{Código modularizado, testes, documentação, Git, GitHub, GitLab, branching, merge requests, revisão de código e processos de implantação.}
\skill{Agilidade e gestão}{Agile, Scrum, Kanban e OKR; planejamento, priorização, refinamento do backlog, reuniões diárias, revisões, retrospectivas e acompanhamento de entregas.}
\skill{Operação e sustentação}{Acompanhamento de produtos de dados novos e legados, suporte pós-implantação, monitoramento de pipelines e modelos, alertas, diagnóstico de falhas e Zabbix.}
\skill{Práticas profissionais}{Qualidade de dados, implantação, monitoramento, documentação técnica, treinamento e comunicação com stakeholders.}
\end{multicols}

\section{Idiomas}

Português · Inglês · Espanhol

\vfill
{\footnotesize\color{grey} Currículo em construção · Última atualização: julho de 2026}

\end{document}
